% GENERATED FILE. Edit canonical lessons, then run npm run generate:books.
% canonical-source-hash: fnv1a64:32da7c0f618dacc4
% canonical-lessons: SA-C307-viyogah, SA-C307-asru, SA-C307-milati, SA-C307-anathah, SA-C307-ksemah

\chapter{Lost and Found}
\label{ch:sa-lost-and-found}
\hlchaptermodality{307}

\begin{chapteropening}
\textbf{By the end of this chapter:} \emph{I can tell how a family was separated and found again.}

Complete the last lesson of chapter 307: I can tell how a family was separated and found again.
\end{chapteropening}

\section[\texorpdfstring{\sk{वियोगः} (viyogaḥ)}{वियोगः (viyogaḥ)}]{\sk{वियोगः} (viyogaḥ) — a separation}

\label{lesson:SA-C307-viyogah}



\begin{quote}
Before the new one: say the Sanskrit for a refugee, then the Sanskrit for comforts.
\end{quote}

\subsection*{You'll want to know: \sk{वियोगः}}
\textbf{\sk{वियोगः}} — \emph{viyogaḥ} — \textquotedblleft{}a separation\textquotedblright{}.

\begin{quote}
\sk{जलौघे} \sk{कुटुम्बस्य} \sk{वियोगः} \sk{अभवत्।}
\end{quote}

\begin{quote}
\emph{jalaughe kuṭumbasya viyogaḥ abhavat.}
\end{quote}

\begin{quote}
In the flood the family was separated.
\end{quote}

A masculine noun.

\begin{cousinweb}
\textbf{\sk{वियोगः}} (\emph{viyogaḥ}) is built from \emph{vi}, \textquotedblleft{}apart\textquotedblright{}, and the root √yuj, \textquotedblleft{}to join\textquotedblright{}; MW gives \emph{viyoga} as disjunction, separation, and the loss of something (MW).
\end{cousinweb}

\begin{grammarlens}[title={What you've built}]
One more word for telling what happened.
\end{grammarlens}

\subsection*{Your turn}
\begin{itemize}
\raggedright
  \item \emph{Say it:} \emph{viyogaḥ}
  \item \emph{Say it:} \emph{viyogaḥ}, once more
  \item \emph{Say it:} \emph{jalaughe kuṭumbasya viyogaḥ abhavat.} — In the flood the family was separated.
  \item \emph{Recall:} say the Sanskrit for a refugee, then the Sanskrit for comforts, then say \emph{viyogaḥ} again
\end{itemize}

\begin{tcolorbox}[breakable,colback=teal!4,colframe=teal!35!black,title={Before you move on}]
What does \sk{वियोगः} mean? (\textquotedblleft{}A separation\textquotedblright{}.) Say it once more.
\end{tcolorbox}
\section[\texorpdfstring{\sk{अश्रु} (aśru)}{अश्रु (aśru)}]{\sk{अश्रु} (aśru) — a tear}

\label{lesson:SA-C307-asru}



\begin{quote}
Before the new one: say the Sanskrit for comforts, then the Sanskrit for a separation.
\end{quote}

\subsection*{You'll want to know: \sk{अश्रु}}
\textbf{\sk{अश्रु}} — \emph{aśru} — \textquotedblleft{}a tear\textquotedblright{}.

\begin{quote}
\sk{मातुः} \sk{नेत्रात्} \sk{अश्रु} \sk{पतति।}
\end{quote}

\begin{quote}
\emph{mātuḥ netrāt aśru patati.}
\end{quote}

\begin{quote}
A tear falls from the mother's eye.
\end{quote}

A neuter noun.

\begin{cousinweb}
\textbf{\sk{अश्रु}} (\emph{aśru}) MW reports that it is supposed to stand for \emph{daśru}, from the root √daṃś, and MW compares Greek \emph{dakru}, Latin \emph{lacryma} and English \emph{tear}; that origin is not settled (MW).
\end{cousinweb}

\begin{grammarlens}[title={What you've built}]
One more word for telling what happened.
\end{grammarlens}

\subsection*{Your turn}
\begin{itemize}
\raggedright
  \item \emph{Say it:} \emph{aśru}
  \item \emph{Say it:} \emph{aśru}, once more
  \item \emph{Say it:} \emph{mātuḥ netrāt aśru patati.} — A tear falls from the mother's eye.
  \item \emph{Recall:} say the Sanskrit for comforts, then the Sanskrit for a separation, then say \emph{aśru} again
\end{itemize}

\begin{tcolorbox}[breakable,colback=teal!4,colframe=teal!35!black,title={Before you move on}]
What does \sk{अश्रु} mean? (\textquotedblleft{}A tear\textquotedblright{}.) Say it once more.
\end{tcolorbox}
\section[\texorpdfstring{\sk{मिलति} (milati)}{मिलति (milati)}]{\sk{मिलति} (milati) — meets; finds again}

\label{lesson:SA-C307-milati}



\begin{quote}
Before the new one: say the Sanskrit for a separation, then the Sanskrit for a tear.
\end{quote}

\subsection*{You'll want to know: \sk{मिलति}}
\textbf{\sk{मिलति}} — \emph{milati} — \textquotedblleft{}meets; finds again\textquotedblright{}.

\begin{quote}
\sk{माता} \sk{शिबिरे} \sk{बालेन} \sk{सह} \sk{मिलति।}
\end{quote}

\begin{quote}
\emph{mātā śibire bālena saha milati.}
\end{quote}

\begin{quote}
The mother meets her child in the camp.
\end{quote}

Past (she …): \sk{सा} \sk{बालेन} \sk{सह} \sk{मिलितवती}, \emph{sā bālena saha militavatī}.

Future: \sk{मिलिष्यति}, \emph{miliṣyati}.

\begin{cousinweb}
\textbf{\sk{मिलति}} (\emph{milati}) comes from the root √mil, \textquotedblleft{}to meet, come together\textquotedblright{} (MW).
\end{cousinweb}

\begin{grammarlens}[title={What you've built}]
One more word for telling what happened.
\end{grammarlens}

\subsection*{Your turn}
\begin{itemize}
\raggedright
  \item \emph{Say it:} \emph{milati}
  \item \emph{Say it:} \emph{milati}, once more
  \item \emph{Say it:} \emph{mātā śibire bālena saha milati.} — The mother meets her child in the camp.
  \item \emph{Recall:} say the Sanskrit for a separation, then the Sanskrit for a tear, then say \emph{milati} again
\end{itemize}

\begin{tcolorbox}[breakable,colback=teal!4,colframe=teal!35!black,title={Before you move on}]
What does \sk{मिलति} mean? (\textquotedblleft{}Meets; finds again\textquotedblright{}.) Say it once more.
\end{tcolorbox}
\section[\texorpdfstring{\sk{अनाथः} (anāthaḥ)}{अनाथः (anāthaḥ)}]{\sk{अनाथः} (anāthaḥ) — orphaned; without a protector}

\label{lesson:SA-C307-anathah}



\begin{quote}
Before the new one: say the Sanskrit for a tear, then the Sanskrit for meets.
\end{quote}

\subsection*{You'll want to know: \sk{अनाथः}}
\textbf{\sk{अनाथः}} — \emph{anāthaḥ} — \textquotedblleft{}orphaned; without a protector\textquotedblright{}.

\begin{quote}
\sk{अनाथः} \sk{बालः} \sk{शिबिरे} \sk{वसति।}
\end{quote}

\begin{quote}
\emph{anāthaḥ bālaḥ śibire vasati.}
\end{quote}

\begin{quote}
The orphaned boy lives in the camp.
\end{quote}

\begin{cousinweb}
\textbf{\sk{अनाथः}} (\emph{anāthaḥ}) is built from \emph{a}, \textquotedblleft{}not\textquotedblright{}, and \emph{nātha}, \textquotedblleft{}a protector, lord\textquotedblright{}; MW gives \emph{anātha} as having no master or protector (MW).
\end{cousinweb}

\begin{grammarlens}[title={What you've built}]
One more word for telling what happened.
\end{grammarlens}

\subsection*{Your turn}
\begin{itemize}
\raggedright
  \item \emph{Say it:} \emph{anāthaḥ}
  \item \emph{Say it:} \emph{anāthaḥ}, once more
  \item \emph{Say it:} \emph{anāthaḥ bālaḥ śibire vasati.} — The orphaned boy lives in the camp.
  \item \emph{Recall:} say the Sanskrit for a tear, then the Sanskrit for meets, then say \emph{anāthaḥ} again
\end{itemize}

\begin{tcolorbox}[breakable,colback=teal!4,colframe=teal!35!black,title={Before you move on}]
What does \sk{अनाथः} mean? (\textquotedblleft{}Orphaned; without a protector\textquotedblright{}.) Say it once more.
\end{tcolorbox}
\section[\texorpdfstring{\sk{क्षेमः} (kṣemaḥ)}{क्षेमः (kṣemaḥ)}]{\sk{क्षेमः} (kṣemaḥ) — safety; well-being}

\label{lesson:SA-C307-ksemah}



\begin{quote}
Before the new one: say the Sanskrit for meets, then the Sanskrit for orphaned.
\end{quote}

\subsection*{You'll want to know: \sk{क्षेमः}}
\textbf{\sk{क्षेमः}} — \emph{kṣemaḥ} — \textquotedblleft{}safety; well-being\textquotedblright{}.

\begin{quote}
\sk{अद्य} \sk{सर्वे} \sk{क्षेमेण} \sk{गृहं} \sk{प्रत्यागच्छन्ति।}
\end{quote}

\begin{quote}
\emph{adya sarve kṣemeṇa gṛhaṃ pratyāgacchanti.}
\end{quote}

\begin{quote}
Today everyone comes home in safety.
\end{quote}

A masculine noun.

\begin{cousinweb}
\textbf{\sk{क्षेमः}} (\emph{kṣemaḥ}) comes from the root √kṣi, \textquotedblleft{}to dwell, abide\textquotedblright{}; MW gives \emph{kṣema} as giving rest or ease and as safety, security (MW).
\end{cousinweb}

\begin{grammarlens}[title={What you've built}]
One more word for telling what happened.
\end{grammarlens}

\subsection*{Your turn}
\begin{itemize}
\raggedright
  \item \emph{Say it:} \emph{kṣemaḥ}
  \item \emph{Say it:} \emph{kṣemaḥ}, once more
  \item \emph{Say it:} \emph{adya sarve kṣemeṇa gṛhaṃ pratyāgacchanti.} — Today everyone comes home in safety.
  \item \emph{Recall:} say the Sanskrit for meets, then the Sanskrit for orphaned, then say \emph{kṣemaḥ} again
\end{itemize}

\begin{tcolorbox}[breakable,colback=teal!4,colframe=teal!35!black,title={Before you move on}]
What does \sk{क्षेमः} mean? (\textquotedblleft{}Safety; well-being\textquotedblright{}.) Say it once more.
\end{tcolorbox}
